% =============================================================================
% System
% =============================================================================

\section{System Design}
\label{sec:system}

\sys{} is a drop-in proxy that sits between Claude Code (CC) and a vLLM
backend serving Llama-3.3-70B-Instruct fp8 with LMCache's
\texttt{LMCacheConnectorV1} attached for cross-request KV reuse.
Figure~\ref{fig:architecture} shows the data path. CC issues
Anthropic-shaped requests against the proxy's \texttt{/v1/messages}
endpoint; the proxy translates each request to vLLM's OpenAI-compatible
\texttt{/v1/chat/completions} API, applies three outbound
transformations, and forwards. The cacheblend integration we present
in this section consists of those three transformations
(\S\ref{sec:cc-segment}--\S\ref{sec:header-norm}); an offline
statistical-span miner (\S\ref{sec:miner}) ships with the artifact
but is not exercised in this paper's measurements.

\begin{figure*}[t]
  \centering
  \begin{tikzpicture}[
    every node/.style={font=\small},
    component/.style={
      draw, rectangle, rounded corners=2pt, minimum height=0.8cm,
      align=center, font=\small\sffamily,
    },
    proxy/.style={component, fill=blue!8, draw=blue!50},
    patch/.style={component, fill=orange!12, draw=orange!60, minimum width=2.3cm, minimum height=0.65cm, font=\footnotesize\sffamily},
    backend/.style={component, fill=green!8, draw=green!55!black},
    sink/.style={component, fill=gray!10, draw=gray!50, font=\footnotesize\sffamily},
    miner/.style={component, fill=purple!8, draw=purple!50, font=\footnotesize\sffamily},
    arrow/.style={-{Stealth[length=2mm]}, thick, draw=black!70},
    offlinearrow/.style={-{Stealth[length=2mm]}, semithick, draw=purple!60, dashed},
  ]

  % CC client (left)
  \node[component, fill=yellow!10, draw=yellow!50!brown, minimum width=2.2cm] (cc) at (0,0)
    {Claude Code};

  % Proxy box (container; title sits INSIDE the top)
  \node[proxy, minimum width=5.4cm, minimum height=3.0cm, anchor=west] (proxybox)
    at ($(cc.east) + (1.2,0)$) {};
  \node[anchor=north, font=\small\sffamily\bfseries, yshift=-0.12cm]
    at (proxybox.north) {\sys{} proxy};

  % Three transforms inside the proxy, vertically stacked, below the title
  \node[patch, anchor=north, yshift=-0.65cm] (p1) at (proxybox.north)
    {(1) CC-segment parser};
  \node[patch, below=0.10cm of p1] (p2) {(2) Header normalization};
  \node[patch, below=0.10cm of p2] (p3) {(3) Pre-seed dispatcher};

  % Backend (right of proxy)
  \node[backend, minimum width=3.6cm, minimum height=3.0cm, anchor=west]
    (backendbox) at ($(proxybox.east) + (1.4,0)$) {};
  \node[anchor=north, font=\small\sffamily\bfseries, yshift=-0.12cm]
    at (backendbox.north) {vLLM + LMCache};
  \node[component, fill=green!15, draw=green!55!black, anchor=north,
        yshift=-0.65cm, minimum width=3.0cm, font=\footnotesize\sffamily]
    (vllm) at (backendbox.north) {vLLM v0.19};
  \node[component, fill=green!20, draw=green!55!black,
        below=0.18cm of vllm, minimum width=3.0cm, minimum height=0.8cm,
        font=\footnotesize\sffamily]
    (lmcache) {LMCache 0.4.2\\\textbf{+ Patch 6}};

  % Offline miner row: placed below the proxy so its arrow
  % into p3 enters the proxy from below, avoiding the title.
  \node[miner, below=1.0cm of proxybox.south, anchor=north, minimum width=4.5cm]
    (minernode) {Offline statistical span miner\\\footnotesize suffix array + Kasai LCP + MinHash @ $J{=}0.8$};
  \node[sink, left=0.5cm of minernode, minimum width=2.0cm, font=\footnotesize\sffamily]
    (corpus) {Training corpus};

  % --- Arrows ---
  % Forward: CC -> proxy -> backend
  \draw[arrow] (cc.east) -- node[above, font=\footnotesize] {\texttt{/v1/messages}} (proxybox.west);
  \draw[arrow] (proxybox.east) -- node[above, font=\footnotesize] {rewritten + chunked} (backendbox.west);

  % Offline: corpus -> miner -> pre-seed dispatcher.
  % Arrow enters proxy from BELOW, terminating at p3 (the pre-seed dispatcher),
  % so it never crosses the proxy title or p1/p2.
  \draw[arrow] (corpus.east) -- (minernode.west);
  \draw[offlinearrow] (minernode.north) -- node[right, font=\footnotesize\itshape, xshift=2pt] {JSONL spans} (p3.south);

  % Pre-seed warmup STOREs travel into lmcache via the normal proxy->backend
  % data path at startup (no extra arrow needed; called out by the {(3)}
  % dispatcher block).

  \end{tikzpicture}
  \caption{\sys{} architecture. The proxy applies three transformations
  to every outbound request: a CC-segment parser (\S\ref{sec:cc-segment})
  that injects cacheblend's separator around recognized structural
  blocks; per-turn header normalization (\S\ref{sec:header-norm}) that
  stabilizes the first KV chunk's hash across turns; and a pre-seed
  dispatcher driven by an offline statistical span miner
  (\S\ref{sec:miner}) that warms the lmcache store with mined repeated
  sequences at startup. Patch~6 (\S\ref{sec:patch6}) is applied to
  lmcache inside the backend container.}
  \label{fig:architecture}
  \Description{Horizontal flow diagram. A Claude Code client on the
  left sends a /v1/messages request to a skillcacher proxy in the
  middle. The proxy contains three stacked components: CC-segment
  parser, header normalization, and pre-seed dispatcher. The proxy
  forwards the rewritten and chunked request to a vLLM-plus-LMCache
  backend on the right; LMCache has Patch 6 applied. Below the proxy,
  an offline statistical span miner reads from a training corpus and
  sends mined JSONL spans up into the pre-seed dispatcher.}
\end{figure*}

\subsection{CC-segment parser}
\label{sec:cc-segment}

\paragraph{The chunking problem.} Cacheblend's segment detector
expects its configured \texttt{blend\_special\_str} (default
\texttt{\#\,\#}) to appear at chunk boundaries in the prompt. Chunks
between separators are hashed and stored individually; chunks below
\texttt{blend\_min\_tokens} (default 256) are ignored. CC's outbound
traffic uses \texttt{\textbackslash n\textbackslash n} as its
predominant block delimiter, which produces 25--180-token segments on
the Llama tokenizer --- below the 256-token floor. Against natural CC
traffic with the published default, the segment detector found zero
chunks and cacheblend's retrieval rate was 0\% regardless of the
workload's textual repetition.

\paragraph{Structural anchors.} The parser identifies semantically
meaningful CC block boundaries by literal anchor matching, then
injects the cacheblend separator on either side of each recognized
block. Anchors are deliberately content-stable across CC versions ---
they key on substrings like \texttt{Status:\textbackslash n} (gitStatus
preamble), \texttt{Recent commits:\textbackslash n}, the
\texttt{<system-reminder>}/\texttt{</system-reminder>} wrapper,
\texttt{Summary:\textbackslash n1.\ } (the 9-section auto-compact
summary), and tool-output markers --- rather than on version digits
or content hashes that drift between releases.

Table~\ref{tab:cc-anchors} lists the recognized block types, their
anchors, and typical token lengths on Llama-3 tokenization. The parser
walks each request's text payloads once, materializes the set of
matched (start, end, kind) spans, deduplicates overlapping matches by
precedence, and then re-emits the original text with separators
inserted at every span boundary. Spans below \texttt{blend\_min\_tokens}
are rolled into the adjacent larger span rather than getting a
separator pair --- this preserves the 256-token floor without forcing
cacheblend to consider tiny isolated blocks.

\begin{table}[t]
\centering\small
\begin{tabular}{@{}llr@{}}
\toprule
\textbf{Block} & \textbf{Anchor (literal)} & \textbf{Tokens} \\
\midrule
CC system header   & \texttt{cc\_version=$\to$Status:}  & 80--120 \\
gitStatus          & \texttt{Status:\textbackslash n}   & 100--400 \\
Recent commits     & \texttt{Recent commits:\textbackslash n} & 50--200 \\
\texttt{<system-reminder>} & wrapper pair                & 30--200 \\
Compaction preamble & \texttt{This session is being$\ldots$} & $\sim$30 \\
Summary 9-section  & \texttt{Summary:\textbackslash n1.\ }    & 400--800 \\
JSONL backreference & \texttt{If you need specific details$\ldots$} & 30--50 \\
\texttt{<command-name>} & wrapper pair                & 5--30 \\
Tool definitions   & (matched by Plan-2 structural tagger) & 300--1500 \\
\bottomrule
\end{tabular}
\caption{Structural anchors recognized by the CC-segment parser. Token
lengths are on Llama-3 tokenizer.}
\label{tab:cc-anchors}
\end{table}

\paragraph{Why proxy-side, not bench-side?}
The parser is wired into the proxy's request-entry path
(\texttt{proxy/server.py:messages}) so it fires on every outbound
request, not only on the benchmark's replay path. Two reasons. First,
CC rebuilds its request body on each turn from its own session state
--- any pre-injection in the user prompt is silently dropped by the
time the body reaches the proxy. Second, putting the parser at the
proxy means the same code unifies live production use (CC pointed at
the proxy) and bench replay (synthetic requests sent through the proxy
under controlled conditions). The validation we report in
\S\ref{sec:evaluation} runs against the proxy's actual rewrite path,
not a parallel implementation.

\subsection{Patches to LMCache 0.4.2}
\label{sec:patches}

The published lmcache + vllm combination (lmcache 0.4.2 against
vllm 0.19.0 with the V1 connector) ships with two integration issues
that surface only on multi-turn agent workloads. The patches that
follow are surgical: Patch~6 (\S\ref{sec:patch6}) is a 4-line
chunk-alignment adjustment inside lmcache's V1 adapter, and the
header normalizer (\S\ref{sec:header-norm}) is a 3-line proxy hook
that runs two regex substitutions per outbound request. The CC-segment
parser (\S\ref{sec:cc-segment}) sits in the proxy and injects
cacheblend's separator at recognized structural anchors.

\begin{table}[h]
\footnotesize
\centering
\begin{tabular}{lc}
\toprule
Stack configuration on natural CC traffic & Failure mode \\
\midrule
Vanilla lmcache 0.4.2 + vllm 0.19.0 (no patches/proxy) & 0\% by mechanism \\
+ Patch~6 only (chunk-aligned \texttt{LOOKUP})         & 0\% by mechanism \\
+ CC-segment parser only (\texttt{\#\,\#} injection)   & 0\% by mechanism \\
+ Patch~6 + parser (no header norm)                    & near 0\% by mechanism \\
\textbf{Full stack} (patches + parser + header norm)   & measured (\S\ref{sec:evaluation}) \\
\bottomrule
\end{tabular}
\caption{Failure-mode analysis of the three patches. The
``0\% by mechanism'' rows are diagnosed analytically rather than
measured end-to-end: without Patch~6, the STORE/LOOKUP chunk-hash
mismatch makes retrieval rate zero by construction
(\S\ref{sec:patch6}); without the parser, cacheblend's
\texttt{blend\_special\_str} never appears in CC traffic and the
segment detector finds zero chunks (\S\ref{sec:cc-segment});
without header normalization, chunk 0's hash drifts every turn and
the chunk-by-chunk LOOKUP short-circuits on the first miss
(\S\ref{sec:header-norm}). Only the full-stack row is measured
end-to-end. A quantitative ablation that runs each intermediate
configuration as a pod replay is follow-up work.}
\label{tab:patch-failure-modes}
\end{table}

\noindent We submit the two upstream-target changes as a candidate
upstream PR~\cite{lmcache-pr-tbd} and apply them in the
container's boot recipe inside
\texttt{scripts/dev/oneshot\_pod.py}.

\subsubsection{Patch 6: chunk-aligned \texttt{LOOKUP}}
\label{sec:patch6}

\paragraph{The bug.}
LMCache's \texttt{LMCacheConnectorV1} computes chunk hashes for the
STORE and LOOKUP paths from two different token-id sources. On the
\textbf{STORE} side, the connector hashes
\texttt{request.token\_ids}, which an upstream vLLM transform has
already chunk-aligned to
$\lfloor n/256 \rfloor \cdot 256$~tokens
(where 256 is \texttt{chunk\_size}). On the \textbf{LOOKUP} side
(\texttt{get\_num\_\allowbreak{}new\_\allowbreak{}matched\_\allowbreak{}tokens}
inside \texttt{vllm\_v1\_\allowbreak{}adapter.py}), the connector
hashes \texttt{request.\allowbreak{}all\_token\_ids}, the full
$n$-token prompt.

For any prompt not exactly a multiple of 256, STORE and LOOKUP hash
different ranges. For example, on a 289-token prompt, STORE hashes
\texttt{tokens[0:256]} but LOOKUP hashes \texttt{tokens[0:289]}; the
hashes never match. Retrieval rate is 0\% regardless of how the
workload is shaped --- this is not a content problem, it's a chunk
boundary problem on the lookup side.

\paragraph{The fix.}
We truncate the LOOKUP path's token range to chunk-size alignment
before hashing, matching the STORE path's range. The patch is a
one-line range adjustment inside
\texttt{lmcache.\allowbreak{}integration.\allowbreak{}vllm.\allowbreak{}vllm\_v1\_adapter}
applied at boot:

\begin{lstlisting}[basicstyle=\ttfamily\footnotesize, breaklines=true, frame=single, framesep=4pt, breakindent=0pt, language=Python]
# patched lookup token computation
chunk = self.cache_engine.chunk_size  # 256
n = len(request.all_token_ids)
aligned_n = (n // chunk) * chunk
token_ids = list(request.all_token_ids[:aligned_n])
\end{lstlisting}

\paragraph{Empirical impact.}
In single-pod commissioning runs (Llama-70B fp8, single H100), the
patch lifted LMCache hit-token-count from 0 (across 40 consecutive
requests with matching content) to non-zero immediately. With the
full \sys{} stack online (patch + parser + normalization),
\S\ref{sec:evaluation} reports rescue rates that reach 95--99\% at
the steady-state peak of the per-turn distribution.

\subsubsection{Per-turn header normalization}
\label{sec:header-norm}

\paragraph{The bug.} CC embeds an
\texttt{x-anthropic-billing-header} field near the front of every
prompt that carries two per-turn / per-build values:
\texttt{cch=<hex>}, a 16-character per-request CC fingerprint, and
\texttt{cc\_version=2.1.X.<suffix>}, where the \texttt{suffix} differs
between regular turns and CC-internal turns (e.g., the
\texttt{/compact} backend call uses a separate build identifier).
Both fields fall within the first 256 tokens of every request. The
first KV chunk's hash therefore differs on every turn of every
session --- chunk 0 is never reused, and because cacheblend's
chunk-by-chunk retrieval short-circuits on the first miss, no
subsequent chunk gets a chance to hit either.

This issue is invisible if you measure cacheblend on synthetic
single-turn requests (the first turn is always a STORE pass), and
invisible if you measure on permuted-RAG benchmarks (every chunk is
re-positioned, so there is no expectation of chunk-0 stability). It
specifically bites \emph{multi-turn agent traffic with a stable system
prompt}, which is the setting cacheblend's compaction-cliff value
proposition rests on.

\paragraph{The fix.}
Before separator injection, the parser rewrites both fields with
stable \texttt{NORM} placeholders:

\begin{lstlisting}[basicstyle=\ttfamily\footnotesize, breaklines=true, frame=single, framesep=4pt, language=Python]
def _normalize_cc_header(text: str) -> str:
    text = _CCH_VAR.sub("cch=NORM;", text)
    text = _CC_VERSION_VAR.sub(
        "cc_version=NORM;", text)
    return text
\end{lstlisting}

The model treats this billing header as opaque text (the backend
itself does not parse it; it is CC-side metadata stuffed into the
prompt for downstream cost attribution), so normalizing its
variable fields should not change downstream token output at $T{=}0$.
We empirically verify this on our setup: on the no-cache condition,
the model's output on normalized vs unnormalized headers is
bit-identical at $T{=}0$ across all 54 (capture, turn) pairs of the
SWE-V+\texttt{post\_compact} subset. We make no claim about
$T{>}0$ sampling or other backends.

After normalization, chunk 0 hashes identically across all turns of one
session. With patch + parser + normalization combined, cacheblend's
per-turn rescue rate on natural CC traffic lifted from a flat 0\% to
the bimodal distribution characterized in \S\ref{sec:eval-hitrate}
(median 87.5\%, with a 13\% mass of turns reaching 95--99\%
steady-state).

\subsection{Statistical span miner (artifact, unmeasured)}
\label{sec:miner}

The structural parser handles CC-specific anchors that are visible in
the source. A complementary component runs offline against a training
corpus of past traces and identifies long-form repeated token
sequences that no source-level anchor exists for --- examples include
fully-rendered tool definitions, repeated multi-line directives in
prompts that the user has cut-and-pasted across sessions, and
near-duplicate file headers in code blocks. \emph{The miner ships as
part of the \sys{} artifact but is not isolated in this paper's
evaluation: all reported numbers use the structural parser plus
Patches 6 and the header normalizer, with statistical-span pre-seed
disabled.} We describe the miner here for completeness; whether and
how much it adds on top of the structural baseline is an open
ablation we defer to follow-up work.

\paragraph{Algorithm.} The miner builds a suffix array over the
concatenated corpus token stream using prefix-doubling construction
(an early implementation using \texttt{sorted(range(n), key=...)} hung
on a 250\,K-token corpus; the prefix-doubling variant runs in
$O(n \log^2 n)$ time and completes on the full \dataset{} corpus in
tens of seconds of pure-Python compute --- see
Appendix~\ref{app:miner} for the per-step breakdown). Kasai's LCP
construction gives the longest-common-prefix array in linear time.
A monotonic-stack walk over LCP emits all maximal-repeat substrings,
including both outer blocks and nested plateaus --- an earlier
single-pass walk missed the nested case and undercounted spans by
30--40\%. The candidate spans are then deduplicated by MinHash at
Jaccard threshold $0.8$ (64 permutations, shingle width 5,
implemented via hand-rolled BLAKE-2b without external dependencies).

\paragraph{Output.} The miner emits a JSONL of canonical spans with
their occurrence counts, source request IDs, and token sequences. On
the \dataset{} corpus (278\,K tokens, 170 streams), the miner
produces 181 unique spans; the highest-frequency hit is a 1691-token
CC system anchor seen in 26 of 170 requests.

\paragraph{Pre-seed.} A proxy-startup hook
(\texttt{pre\_seed\_statistical\_spans}) reads the JSONL and, for each
span, dispatches a synthetic STORE request whose prompt contains the
span; the response is discarded. After startup, every cached span is
present in the lmcache store and is reachable on its first real
production reference. Pre-seed is gated by the
\texttt{SKILLCACHER\_STATISTICAL\_SPANS\_FILE} environment variable
--- unset for the structural-only baseline, set to the JSONL path
for the structural+statistical condition. The two configurations
share all other code; the only difference is whether the pre-seed
warmup fires.
